\section{Piecewise quadratic certificates in dimension one}
\label{sec:piecewise-quadratic}

In this section, we will study a failure regime of the piecewise
linear certificates in dimension one and see how it leads to sufficiency
of piecewise quadratic certificates.

Additionally, we observe an interesting connection to the \textit{Taylor expansion} of the best possible
certificate, namely a martingale certificate for which the generator is zero.

\subsection{Failure of piecewise linear certificates in dimension one}

We will consider the setup from \autoref{sec:piecewise-linear} again, but with
additional assumptions on absorption at the boundary, an increase condition on
the piecewise linear function $V(x)$, and a non-degeneracy condition on the
diffusion term $g(x)$.

The key point is that the piecewise linear setup is not flexible enough when
the certificate is required to increase in the same direction as the drift. In
that case, a piecewise linear function has no curvature term which could
compensate for the positive drift contribution.

So essentially, we will show that:
\begin{fact}
(\textit{Informally.}) There is a case where no piecewise linear certificate can make $V(X_t)$ a
supermartingale.
\end{fact}

Although it should be clear from previous sections, for completeness, the setup
is as follows:

\begin{setup}
Let:
    \begin{enumerate}
        \item $dX_t = f(X_t)dt + g(X_t)dW_t$ be a one-dimensional SDE.
        \item $V(x)$ be a piecewise linear function such that
        $V(x)=a_i x+b_i$, where $x\in(x_i,x_{i+1})$.
        \item $V(x)$ is continuous at breakpoints:
        $V(x_i^+)=V(x_i^-)=V(x_i)$.
        \item as a consequence, $V'_{-}(x_i) = a_{i-1}$, $V'_{+}(x_i) = a_i$
        \item The $X_t$ process is \textbf{absorbed} at the boundary of
        $I=[a,b]$.
        \item $V(x)$ has the \textbf{increase condition} at the boundary:
        \begin{equation}
        V(a)<V(b).
        \end{equation}
        \item $g(x)^2>0$ for all $x\in(a,b)$, i.e. the \textbf{diffusion term is
        non-degenerate}.
        \item $M_t$ is a martingale, i.e. the \textbf{martingale condition}.
        \end{enumerate}
\end{setup}

Motivated by the example given in \autoref{sec:failure-of-piecewise-linear},
we will now show that there is a broad class of SDEs for which increasing
piecewise linear certificates cannot be supermartingales.

\begin{fact}
Assume that for all $\varepsilon>0$, there exists
$x\in(a,a+\varepsilon)$ such that
\begin{equation}
f(x)>0.
\label{eq:positive-drift-arbitrarily-close}
\end{equation}
Then there is no continuous piecewise linear $V$ satisfying the increase
condition
\begin{equation}
V(a)<V(b)
\end{equation}
and the concavity condition
\begin{equation}
V'_{+}(x_i)-V'_{-}(x_i)\leq 0
\end{equation}
at every breakpoint, such that $V(X_t)$ is a supermartingale.

\textbf{Proof}. Assume, for contradiction, that such a piecewise linear function $V$ exists
and that $V(X_t)$ is a supermartingale.

Let $p(x)=V'(x)$ on each linear region. Since $V$ is piecewise linear, the
Itô--Tanaka formula gives
\begin{align}
    V(X_t)
    &=
    V(X_0)
    +
    \int_0^t p(X_s)f(X_s)\,ds                                      \\
    &\quad
    +
    \int_0^t p(X_s)g(X_s)\,dW_s                                    \\
    &\quad
    +
    \frac12\sum_i
    \left(V'_+(x_i)-V'_-(x_i)\right)L_t^{x_i}(X).
\end{align}
The stochastic integral is a martingale by assumption. Therefore, if
$V(X_t)$ is a supermartingale, the finite-variation part cannot create
positive drift on any linear region. Hence the necessary interior condition
is
\begin{equation}
    p(x)f(x)\leq 0
    \qquad
    \text{on every linear region.}
    \label{eq:necessary-interior-condition-pl}
\end{equation}

This also follows from \hyperref[fact:nonpositive-drift-supermartingale]
  {Fact~\ref*{fact:nonpositive-drift-supermartingale}}
in \autoref{sec:piecewise-linear}.

We now use the concavity condition. Since
\begin{equation}
    V'_+(x_i)-V'_-(x_i)\leq 0
\end{equation}
at every breakpoint, the slopes of $V$ are nonincreasing from left to
right. Therefore $V$ is concave.

Let $(a,x_1)$ be the first linear region. We claim that the slope of $V$ on
this first region must be positive. Indeed, suppose instead that the first
slope were nonpositive. Since the slopes are nonincreasing, all later
slopes would also be nonpositive. Thus,
\begin{align}
    V(b)-V(a)
    &=
    \int_a^b V'(x)\,dx                                             \\
    &\leq 0.
\end{align}
This contradicts the increase condition $V(a)<V(b)$. Hence the first slope
must be positive, so
\begin{equation}
    p(x)>0
    \qquad
    \text{for all }x\in(a,x_1).
    \label{eq:first-slope-positive-pl}
\end{equation}

Now use the assumption on the drift. Since for every $\varepsilon>0$ there
exists $x\in(a,a+\varepsilon)$ such that $f(x)>0$, we may choose
$\varepsilon=x_1-a$. Then there exists $x^\ast\in(a,x_1)$ such that
\begin{equation}
    f(x^\ast)>0.
\end{equation}
Combining this with \eqref{eq:first-slope-positive-pl}, we get
\begin{equation}
    p(x^\ast)f(x^\ast)>0.
\end{equation}
This contradicts the necessary interior condition
\eqref{eq:necessary-interior-condition-pl}. Therefore no such piecewise
linear $V$ can make $V(X_t)$ a supermartingale.

\end{fact}

The conclusion is that piecewise linear certificates fail when the certificate
is forced to increase in the same direction as the drift. The obstruction comes
from the fact that piecewise linear functions have no interior curvature term.
Indeed, away from the kinks,
\begin{equation}
\mathcal L V(x)=V'(x)f(x),
\end{equation}
so a positive drift contribution cannot be compensated. This motivates the
piecewise quadratic case, where the additional curvature term
\begin{equation}
\frac12 V''(x)g(x)^2
\end{equation}
can offset the positive drift contribution.

\subsection{Sufficiency of piecewise quadratic certificates in dimension one}

In this subsection we will show that piecewise quadratic certificates are sufficient 
to guarantee that $V(X_t)$ is a supermartingale. 

Let $V$ be piecewise quadratic on $I=[a,b]$, so that on every interval
$(x_i,x_{i+1})$ we have
\begin{equation}
    V(x)=a_ i x^2+ b_i x+ c_i .
\end{equation}
Assume also that $V$ is continuous at the breakpoints. Then
\begin{equation}
    V'(x)=2a_i x+b_i,
    \qquad
    V''(x)=2a_i,
    \qquad x\in(x_i,x_{i+1}).
\end{equation}
Therefore the generator on the interior of a quadratic piece is
\begin{equation}
    \mathcal L V(x)
    =
    \left(2a_i x+b_i\right)f(x)
    +
    a_i g(x)^2.
    \label{eq:piecewise-quadratic-generator}
\end{equation}

Here, there are always exists a solution $(a_i,b_i)$ to the interior condition $\mathcal L V(x)\leq 0$ on each piece.

This gives the following sufficient condition.

\begin{fact}
Let $V$ be continuous and piecewise quadratic on $I=[a,b]$. Suppose that:
\begin{enumerate}
    \item on every open piece $(x_i,x_{i+1})$,
    \begin{equation}
        \mathcal L V(x)
        =
        V'(x)f(x)+\frac12 V''(x)g(x)^2
        \leq 0;
        \label{eq:pq-interior-condition}
    \end{equation}
    \item at every breakpoint,
    \begin{equation}
        V'_+(x_i)-V'_-(x_i)\leq 0;
        \label{eq:pq-kink-condition}
    \end{equation}
    \item the stochastic integral
    \begin{equation}
        M_t=\int_0^t V'(X_s)g(X_s)\,dW_s
    \end{equation}
    is a martingale.
\end{enumerate}
Then $V(X_t)$ is a supermartingale.

\textbf{Proof}. By the Itô--Tanaka formula for a piecewise $C^2$ function,
\begin{align}
    V(X_t)
    &=
    V(X_0)
    +
    \int_0^t
    \left(
        V'(X_s)f(X_s)+\frac12 V''(X_s)g(X_s)^2
    \right)ds                                                   \\
    &\quad+
    \int_0^t V'(X_s)g(X_s)\,dW_s                                \\
    &\quad+
    \frac12\sum_i
    \left(V'_+(x_i)-V'_-(x_i)\right)L_t^{x_i}(X).
\end{align}
The first term is nonincreasing by
\eqref{eq:pq-interior-condition}. The stochastic integral is a martingale by
assumption. Finally, since local time is nondecreasing in $t$, the kink term is
nonincreasing by \eqref{eq:pq-kink-condition}. 
Therefore $V(X_t)$ is a supermartingale.
\end{fact}

\subsection{Connection to Taylor expansion of the best possible certificate}

Let us now again consider the SDE
\begin{equation}
    dX_t=dt+\sqrt 2\,dW_t,
\end{equation}
the generator is
\begin{equation}
    \mathcal L V(x)=V'(x)+V''(x).
\end{equation}
A martingale certificate is obtained by asking for the generator to vanish:
\begin{equation}
    V'(x)+V''(x)=0.
    \label{eq:martingale-certificate-ode}
\end{equation}
Solving this ODE gives
\begin{equation}
    V(x)=C_1+C_2 e^{-x}.
\end{equation}
If we impose the normalization $V(0)=0$ and $V(1)=1$, then
\begin{equation}
    V_\star(x)
    =
    \frac{1-e^{-x}}{1-e^{-1}}.
    \label{eq:best-certificate-exp}
\end{equation}
This is the martingale certificate: $\mathcal L V_\star(x)=0$ on the interior.

Now expand the numerator of \eqref{eq:best-certificate-exp} around $0$:
\begin{equation}
    1-e^{-x}
    =
    x-\frac{x^2}{2}+\frac{x^3}{6}-\cdots .
\end{equation}
The quadratic certificate used above is exactly the second-order truncation of
this numerator:
\begin{equation}
    V_{\mathrm{quad}}(x)
    =
    x-\frac{x^2}{2}.
\end{equation}
So, up to the constant $(1-e^{-1})^{-1}$ and up to truncation
error, the quadratic certificate is the Taylor approximation of the martingale
certificate.
